\begin{table}[ht]
    \centering
    \caption{Perbandingan strategi retrieval.}
    \label{tab:bab2.3-strategi-retrieval}

    \begin{tabular}{p{3.2cm} p{4.0cm} p{3.4cm} p{3.4cm}}
        \hline
        \textbf{Strategi}
        & \textbf{Dasar pencocokan}
        & \textbf{Kelebihan}
        & \textbf{Keterbatasan} \\
        \hline

        Vector retrieval
        &
        Kemiripan representasi semantik antara kueri dan dokumen.
        &
        Dapat menangani variasi istilah dan parafrasa.
        &
        Bergantung pada kualitas representasi dan kesesuaian domain.
        \\

        Lexical retrieval (BM25)
        &
        Kecocokan istilah pada kueri dan dokumen dengan pembobotan term.
        &
        Hubungan antara istilah kueri dan dokumen lebih langsung.
        &
        Lebih sensitif terhadap perbedaan istilah dan parafrasa.
        \\

        Hybrid retrieval
        &
        Gabungan sinyal semantik dan leksikal.
        &
        Dapat memanfaatkan kedua bentuk kecocokan.
        &
        Menambah kompleksitas dan tetap bergantung pada kualitas kandidat awal.
        \\

        \hline
    \end{tabular}
    
    \noindent\tabnotestyle Sumber: Disusun penulis berdasarkan \textcite{manning2008} dan \textcite{xiong2024}.\normalsize

\end{table}